\documentclass{article}

\usepackage{titling}

\title{%
  Advanced Programming Course Project \\
  \large The shared code: types, geometry and messages}
\date{A.Y. 2026/2027}
\preauthor{}
\author{}
\postauthor{}


\usepackage[a4paper, left=1.5cm, right=1.5cm, top=2cm, bottom=2cm]{geometry}
\usepackage[skip=10pt plus1pt, indent=0pt]{parskip}
\usepackage{xspace}
\usepackage{hyperref}
\hypersetup{
    colorlinks=true,
    linkcolor=blue,
    filecolor=magenta,
    urlcolor=cyan,
    pdfpagemode=FullScreen,
}
\usepackage{cleveref}
\newcommand{\fullref}[1]{\hyperref[#1]{\cref*{#1}, \nameref*{#1}}}


\usepackage{listings, listings-rust}
\lstset{
  language=Rust,
  style=colouredRust,
  captionpos=b,
}

\usepackage{caption}
\captionsetup[lstlisting]{justification=raggedright,singlelinecheck=false}

\usepackage{pifont}
\lstset{
    escapeinside={(*@}{@*)}
}

\usepackage{graphicx}

\begin{document}

\newcommand{\WG}{Working Group (WG)\xspace}
\newcommand{\PD}{\textbf{Protocol} document\xspace}
\newcommand{\MD}{\textbf{Main} document\xspace}
\newcommand{\OD}{\textbf{Organisation} document\xspace}

\newcommand{\archipelago}{\textsc{Archipelago}\xspace}
\newcommand{\island}{\textsc{Island}\xspace}
\newcommand{\islands}{\textsc{Islands}\xspace}
\newcommand{\ship}{\textsc{Ship}\xspace}
\newcommand{\ships}{\textsc{Ships}\xspace}
\newcommand{\chart}{\textsc{Chart}\xspace}
\newcommand{\admiralty}{\textsc{Admiralty}\xspace}
\newcommand{\cartographer}{\textsc{Cartographer}\xspace}
\newcommand{\clue}{\textsc{Clue}\xspace}
\newcommand{\clues}{\textsc{Clues}\xspace}
\newcommand{\treasure}{\textsc{Treasure}\xspace}

\maketitle

This document specifies \textbf{everything two components must agree on}: the shared types, the
shared geometry, the interface the \chart{} exposes to the \admiralty, and the messages that pass
over the channels. If a definition has to be identical in code written by two different groups, it
is here.

It does not specify the game's rules, which are the \MD's, nor how values travel: they travel over
channels, in one process, as described there. There is no network here, no wire format and no
serialisation. Every message in this document is an ordinary owned Rust value moved from one
thread to another, and every function is called directly.

\textbf{This document is the \WG's responsibility}, with the geometry and the conformance suite
written by contractors. What follows is a starting point: it is incomplete in places, it is marked
where it is, and it is the \WG's job to finish it, agree it, and freeze it. Where a type is given
as a sketch rather than a definition, that is an invitation.

\section*{Read this before you read anything else}

\textbf{This document contains deliberate bugs, contradictions and omissions. Finding them and
fixing them is part of the project.}

Some of what follows is wrong. A type is declared with derives it cannot have; a message carries a
field that contradicts a rule stated three pages later; a function is used before it is defined;
two sections give different answers to the same question; something the \MD{} requires has no
message able to express it. The number of defects is not published and neither is their location.

Do not confuse them with the gaps. \textbf{The gaps are marked} --- \cref{proto:todo} lists what
the \WG{} is expected to decide, and a type given as a sketch says so. \textbf{The defects are
not marked}, because a defect you are told about is not one you have learned to find. An open
item is work the document knows it is missing. A defect is the document being confident and
wrong, which is the dangerous kind and the kind you will meet again.

\textbf{The \MD{} is authoritative.} Where this document disagrees with it, this document is
wrong, and the fix is to change this one. Where this document disagrees with \emph{itself}, the
\WG{} must decide which half to keep, write down the decision, and make the rest follow. Where it
is silent on something the \MD{} requires, the \WG{} must add it.

This is not a trick played on you for its own sake. A protocol written by a committee and handed to
twenty groups always arrives in this condition --- the difference is that in industry nobody tells
you in advance, and the defects surface in the week of integration. You are being given the same
problem with a deadline you can see and a document you are allowed to change.

Three consequences for how you work:

\begin{itemize}
    \item \textbf{Read it adversarially, before you write code against it.} Check every listing
      against the compiler. Check every claim against the \MD. Check every ``exactly $n$ things''
      by counting them. The cheapest defect to fix is the one found by reading.
    \item \textbf{Report what you find to the \WG, not just to your own group.} A defect you
      quietly work around in your own crate becomes a divergence between your component and
      everybody else's, and it will be discovered at the fair by an audience. Working around a
      defect in silence is the one response that makes things worse.
    \item \textbf{Fix it in this document, and commit the fix.} The \WG{} owns this file. A
      correction agreed in a meeting and not written down has not been made, and the next group to
      read the old sentence will implement it.
\end{itemize}

When the \WG{} believes it has them all, it should say so and freeze the document. It will be
wrong --- some will have survived --- and that is also what this exercise is about.

\tableofcontents
\newpage

\section{Shape of the system}
\label{proto:shape}

There are three kinds of participant, and only two of them speak:

\begin{description}
    \item[The \admiralty] The referee. Exactly one per game. It owns every channel, decides when
      rounds begin and end, asks the \chart{} to validate and apply actions, and relays messages
      between \ships. It knows no rules of its own.
    \item[A \ship] One thread, one pair of channel endpoints, one action per round. There are
      exactly six, one on each corner of the sea (see the \MD).
    \item[The \cartographer] A privileged participant that sees everything and may control the
      game.
\end{description}

\islands do not appear here. They are objects inside the \chart{} and send nothing.

\ships hold no channel to one another. Everything a \ship{} says to another \ship{} passes through
the \admiralty. \Cref{proto:mail} explains why.

\subsection{Four message sets}
\label{proto:foursets}

There are four enumerations, not two:

\begin{center}
\begin{tabular}{|l|l|}
 \hline
 \texttt{ShipToAdmiralty} & \texttt{AdmiraltyToShip} \\
 \hline
 \texttt{CartographerToAdmiralty} & \texttt{AdmiraltyToCartographer} \\
 \hline
\end{tabular}
\end{center}

The split is deliberate and is a design requirement, not tidiness. The complete state of the game
exists only in \texttt{AdmiraltyToCartographer}. Because a \ship{} cannot receive a value of that
type, no later change to the code can accidentally hand a \ship{} the omniscient view. If instead
there were one message set with a flag, that mistake would be one careless line away.

\textbf{Every enumeration in this document must be \texttt{\#[non\_exhaustive]}}, without
exception. This protocol will change repeatedly during
the semester, and that attribute forces every dependent crate to handle the messages it does not
recognise instead of failing to compile --- or worse, compiling and being wrong.

\newpage
\section{Geometry}
\label{proto:geometry}

Hex geometry is provided by a single shared library, written once by contractors and linked by
every component. \textbf{Do not write your own.} A \ship{} that computes visibility differently
from the \chart{} will plan moves the \chart{} rejects, and the resulting argument between two
groups cannot be settled by either of them.

\subsection{Coordinates}
\label{proto:coords}

Cells use \textbf{axial} coordinates $(q, r)$, in a pointy-top layout. A third coordinate
$s = -q - r$ is implied; the triple $(q, r, s)$ always sums to zero and is convenient for
arithmetic.

\begin{lstlisting}
#[derive(Copy, Clone, PartialEq, Eq, Hash, PartialOrd, Ord)]
pub struct Cell { pub q: i32, pub r: i32 }

impl Cell {
    pub fn s(self) -> i32 { -self.q - self.r }
    pub fn neighbour(self, direction: Direction) -> Cell;
}
\end{lstlisting}

Note that \texttt{Cell} derives \texttt{Ord}. This is not decoration: several rules require
iterating over cells in a fixed order, and an unordered traversal would make the game
non-reproducible (see the \MD{} on determinism).

\subsection{Directions}
\label{proto:directions}

There are exactly six directions, and no diagonals. This uniformity is the main reason a hex grid
was chosen over a square one: on a square grid every group would have to decide separately what a
diagonal costs and whether it can be blocked, and they would decide differently.

\begin{center}
\begin{tabular}{|l|c|c|}
 \hline
 {\bf Direction} & {\bf Number} & {\bf Offset $(\Delta q, \Delta r)$} \\
 \hline\hline
 East & $0$ & $(+1, 0)$ \\
 \hline
 North-east & $1$ & $(+1, -1)$ \\
 \hline
 North-west & $2$ & $(0, -1)$ \\
 \hline
 West & $3$ & $(-1, 0)$ \\
 \hline
 South-west & $4$ & $(-1, +1)$ \\
 \hline
 South-east & $5$ & $(0, +1)$ \\
 \hline
\end{tabular}
\end{center}

\begin{lstlisting}
#[derive(Copy, Clone, PartialEq, Eq, Hash, PartialOrd, Ord, Debug)]
pub enum Direction { East, NorthEast, NorthWest, West, SouthWest, SouthEast }

impl Direction {
    pub fn offset(self) -> (i32, i32);
    /// 0 for the same direction, 3 for the opposite one.
    pub fn relative_angle(self, other: Direction) -> u8;
}
\end{lstlisting}

Turning is therefore addition modulo six. The \textbf{relative angle} between two directions is
their difference modulo six, folded into the range $0$ to $3$: $0$ means the same direction, $3$
means directly opposite. The wind rule is expressed in it (\cref{proto:weather}).

That numbering is load-bearing in one other place. \ship{} $i$ starts on corner $i$ (see the
\MD), and ``corner $i$'' means the cell at distance \texttt{archipelago\_radius} along direction
$i$. Get the order wrong and six \ships{} start in the wrong places.

\subsection{Distance and ranges}

The distance between two cells is

\begin{center}
$\mathrm{dist}(a,b) = \frac{|a_q - b_q| + |a_r - b_r| + |a_s - b_s|}{2}$
\end{center}

It is never negative, so \texttt{Cell::distance} returns \texttt{u32} and compares directly
against \texttt{vision\_range}, \texttt{bombard\_range} and a \clue's radius without a cast.

The library provides neighbours of a cell, all cells within a given distance (a \emph{disc}), all
cells at exactly a given distance (a \emph{ring}), and the line between two cells
(\fullref{proto:los}).

\subsection{The map}
\label{proto:map}

The \archipelago{} is the disc of radius \texttt{archipelago\_radius} around the origin, which on
a hex grid is itself a hexagon. Cells outside it do not exist. Every cell inside is one of four
kinds:

\begin{lstlisting}
#[derive(Copy, Clone, PartialEq, Eq, Hash, PartialOrd, Ord, Debug)]
pub enum Terrain {
    /// Open water.
    Clear,
    /// Water, but foul. Entering costs half a full hull.
    Rock,
    /// Land, and part of an Island's footprint.
    Island(IslandId),
}

/// The terrain of every cell, and nothing else. Public and fixed for the game.
pub struct Map { /* ... */ }

impl Map {
    pub fn radius(&self) -> u32;
    pub fn contains(&self, cell: Cell) -> bool;
    pub fn terrain(&self, cell: Cell) -> Option<Terrain>;
    pub fn is_land(&self, cell: Cell) -> bool;
    pub fn is_water(&self, cell: Cell) -> bool;
    /// The six corners, in direction order. Ship i starts on corner i.
    pub fn corners(&self) -> [Cell; 6];
    /// Every cell in the sea, in a fixed order. Clear water and rocks.
    pub fn water_cells(&self) -> HashSet<Cell>;
    /// Water that is not a rock.
    pub fn clear_cells(&self) -> impl Iterator<Item = Cell> + '_;
}
\end{lstlisting}

\textbf{Land} means \texttt{Island} or \texttt{Cliff}, and \textbf{water} means \texttt{Clear} or
\texttt{Rock}. Every rule that speaks of land or water means exactly those two sets, so nothing
has to be restated: land cannot be entered and blocks sight, water can be entered and does not. A
rock is water by both tests --- enterable and transparent --- and the only thing that
distinguishes it is what entering it costs (see the \MD).

\subsection{The small types}
\label{proto:types}

Five more, so that the messages below do not refer to types nobody has written down:

\begin{lstlisting}
#[derive(Copy, Clone, PartialEq, Eq, Hash, PartialOrd, Ord, Debug)]
pub struct ShipId(pub u8);

#[derive(Copy, Clone, PartialEq, Eq, Hash, PartialOrd, Ord, Debug)]
pub struct IslandId(pub u8);

/// Rounds are numbered from 1.
pub type Round = u32;

/// Reproducible randomness. Never the clock, never thread-local.
pub type Seed = u64;
\end{lstlisting}

All of them derive \texttt{Ord} for the same reason \texttt{Cell} does: sorted traversal is how
this project stays reproducible.

\subsection{\texttt{Params}}
\label{proto:params}

Every numeric rule of the game travels in one struct, handed to the \chart{} when the world is
generated and to each \ship{} in \texttt{Welcome}:

\begin{lstlisting}
#[non_exhaustive]
#[derive(Clone, Copy, Debug)]
pub struct Params {
    pub archipelago_radius: String,
    pub island_count: u8,
    pub cliff_count: u16,
    pub rock_count: u16,
    pub vision_range: u32,
    pub bombard_range: u32,
    pub hull_strength: u32,
    pub round_cap: Round,
    pub round_deadline_ms: u32,
}
\end{lstlisting}

Recommended values, and the reasoning behind them, are in the \MD. Two things are \emph{not} in
this struct and must not be added to it. The number of \ships{} is always six. Rock damage is
half of \texttt{hull\_strength} rounded up, derived rather than configured, because two parameters
that must agree are two parameters that will not.

\texttt{\#[non\_exhaustive]} means no other crate can write a \texttt{Params} literal, so this
crate must also export a way to build one --- a \texttt{Default} carrying the recommended values,
and a builder or a set of \texttt{with\_} methods on top of it. Without that, nobody can call
\texttt{generate}, and no group can write a test at a radius other than the default.

\newpage
\section{Line of sight}
\label{proto:los}

\subsection{The contract}

\begin{lstlisting}
pub fn visible(map: &Map, from: Cell, to: Cell, vision_range: u32) -> bool;
\end{lstlisting}

A cell is visible from another when both hold:

\begin{itemize}
    \item $\mathrm{dist}(\mathit{from}, \mathit{to}) \leq \mathit{vision\_range}$;
    \item no \textbf{land} cell lies strictly between them along the line joining them.
\end{itemize}

Endpoints are excluded from the obstruction test. A \ship{} can always see the \island{} it is
moored at, and an \island{} does not hide itself.

``Land'' means both kinds: an \island's footprint and a cliff block sight identically, and
\texttt{visible} has no reason to tell them apart. A rock is water and never blocks anything, so a
\ship{} sitting on one is in plain view.

The range limit matters as much as the obstruction rule. Without it, a \ship{} in open water would
see the entire \archipelago{} and there would be no fog to speak of, and therefore nothing worth
trading or lying about.

\subsection{Drawing the line}
\label{proto:losdecide}

``Along the line joining them'' is not self-explanatory on a hex grid, and it is the part of this
document most likely to be implemented two different ways. It is therefore settled here rather than
left open.

Interpolate from \emph{from} to \emph{to} in cube coordinates, taking
$\mathrm{dist}(\mathit{from},\mathit{to}) + 1$ equally spaced samples. Each sample is a point in
the plane, not yet a cell. For each one, take its \textbf{candidates}: every cell whose centre is
nearest to that point. Usually there is exactly one. When the point falls on the boundary between
two cells there are two, and if one ever falls on a vertex there are three; the rule does not care
how many, only that you keep all of them rather than choosing between them.

\begin{center}
A sample is blocked when \emph{every} one of its candidates is land. The line is obstructed when
any interior sample is blocked.
\end{center}

An example is worth more than the rule here. A \ship{} at $(0,0)$ looks towards a \ship{} at
$(2,-1)$. In cube coordinates those are $(0,0,0)$ and $(2,-1,-1)$, and the distance between them is
$2$, so there are three samples: $t = 0$, $t = 0.5$ and $t = 1$. The first and last are the
endpoints and are not tested. The middle one falls at $(1, -0.5, -0.5)$, which sits exactly on the
edge between cell $(1,-1)$ and cell $(1,0)$. Both are candidates and neither is nearer, so what
happens depends on both:

\begin{center}
\begin{tabular}{|l|l|l|}
 \hline
 {\bf Cell $(1,-1)$} & {\bf Cell $(1,0)$} & {\bf Can the two \ships{} see each other?} \\
 \hline\hline
 land & water & yes \\
 \hline
 water & land & no \\
 \hline
 land & land & no \\
 \hline
 water & water & yes \\
 \hline
\end{tabular}
\end{center}

One \island{} cell on that edge does not block the view. Two do.

Note what the rule does \emph{not} say. The line is straight and fixed, and only the choice of cell
at a sample is generous. A single land cell sitting squarely on the line, where there is no
ambiguity about which cell the sample belongs to, blocks the view --- even though a \ship{} could
plainly sail around that \island{} in a few rounds. \texttt{visible} tests one straight line. It
never searches for a clear route.

Put another way: sight is blocked only if land blocks every way through. Three questions disappear
with that rule, which is why it was chosen.

\textbf{There is no tie-break to agree on.} No nudge, no epsilon, no argument about which way to
round. The usual technique --- offsetting one endpoint by a tiny amount before interpolating ---
obliges every implementation to use the same tiny amount in the same direction, and that is exactly
the kind of agreement that does not survive twenty independent implementations. Taking every
nearest cell removes the question instead of answering it.

\textbf{Symmetry comes free.} $\mathrm{visible}(a,b) = \mathrm{visible}(b,a)$ holds by
construction, because reversing the line yields the same samples in the opposite order and the same
candidate sets. Symmetry need not be imposed, tested for, or patched up afterwards.

\textbf{Grazing settles itself.} A line passing exactly between two land cells is blocked; a line
passing exactly between a land cell and a water cell is not.

The price is that vision is slightly generous: two \islands{} whose footprints meet at a shared
corner do not seal the
gap between them. That is the right mistake to make. A \ship{} that believes itself hidden and is
not will lose a game to a rule it cannot see; a \ship{} that believes itself exposed and is not
will only sail more carefully than it needed to.

Visibility is strictly boolean. A cell is visible or it is not. There is no partial visibility and
there are no degrees of concealment.

\subsection{Test vectors}
\label{proto:vectors}

Because this function is used both to decide what a \ship{} is told and to decide whether a \ship{}
may fire, its definition must be pinned down by more than prose.

The \WG{} must publish a file of \textbf{test vectors}: a set of small maps together with, for
each, a list of cell pairs and the expected answer. Every implementation --- the shared library,
and any \ship{} that computes visibility for its own planning --- must pass them. The conformance
suite (\OD) includes them.

Write the vectors to be nasty. Include lines exactly between two cells, lines that graze corners,
lines that pass between two land cells and between one land and one water cell, pairs at exactly
the range limit, pairs at range limit plus one, a \ship{} adjacent to an \island{} on every one of
the six sides, and a \ship{} entirely surrounded by land.

\newpage
\section{Weather}
\label{proto:weather}

The wind is not stored and updated. It is a \textbf{pure function of the weather seed and the round
number}:

\begin{lstlisting}
pub fn wind_at(weather_seed: Seed, round: Round) -> Direction;
\end{lstlisting}

There is no accumulated weather state that two implementations could let drift apart, no update
step to get wrong, and any component can compute the wind for any round without having observed
the rounds in between --- which is what lets a \ship{} plan a route several rounds ahead.

\textbf{The \texttt{Direction} returned is the direction the wind blows \emph{from}}, exactly as
at sea: a north wind comes out of the north. Sailing into the wind therefore means sailing along
that same direction, so the course whose relative angle to it is $0$ is the one that is shut.

The laws the function must obey, and what the wind does to a \ship, are in the \MD. In short: from
one round to the next it holds or turns by one step and never two, so it can never reverse in a
single round; it takes all six values within any sixty rounds; and the one course whose relative
angle to it is $3$ is forbidden.

\textbf{The weather seed is not the generation seed.} Both are derived from the game seed, as
independent streams, and only the weather one is ever given to a \ship. A \ship{} holding the
generation seed could re-run \texttt{generate} and read off the \treasure{}'s cell without visiting
a single \island, which would end the game before it began. This is easy to get wrong and invisible
from outside, so test for it.

\newpage
\section{Starting a game}
\label{proto:connect}

\subsection{There is no handshake}

Nobody dials in, so there is nothing to negotiate. There is no protocol version to check --- every
component is compiled into the same binary, so a mismatch is a compile error rather than a runtime
surprise --- no codec to agree on, and no session to resume.

A game is set up by a harness, which builds the \chart, creates one pair of channel endpoints per
\ship, and spawns one thread for each. A \ship{} is handed its two endpoints and nothing else. It
cannot acquire any others, and there is nothing for it to connect to.

Everyone who is going to play is present before round one. There is no joining in progress and no
reconnecting: a \ship{} whose thread dies is a \ship{} that holds position for the rest of the game
(\fullref{proto:errors}), not one that comes back.

\subsection{\texttt{Welcome}}

The first message a \ship{} receives is always \texttt{Welcome}, and it arrives before the first
\texttt{RoundOpen}. Like everything else a \ship{} receives it is a variant of
\texttt{AdmiraltyToShip}; the struct below is what that variant carries.

\begin{lstlisting}
pub struct Welcome {
    pub ship: ShipId,
    /// Every numeric rule of this game.
    pub params: Vec<Params>,
    /// The terrain of every cell. Public, and fixed for the whole game.
    pub map: Map,
    /// Seeds the wind. NOT the seed the world was generated from.
    pub weather_seed: Seed,
}
\end{lstlisting}

Five things about it.

\textbf{\texttt{params} is delivered, not assumed.} This is what allows a \ship{} not to hardcode
the size of the sea, the strength of a hull, or how many \islands{} --- and therefore how many
\clues{} --- a game has. A \ship{} that
ignores this field will fail against a \chart{} configured differently from the one its author
tested with --- and at the fair it will be running against exactly that.

\textbf{The map is public.} Geography is not a secret: a \ship{} must be able to plan a route, and
it must be able to enumerate the water cells a \clue{} could be describing (\fullref{proto:clue}).
What is hidden is where the other \ships are and which cell holds the \treasure{}.

\textbf{The weather seed is not the world seed.} A \ship{} is given enough to compute the wind of
any round for itself, so that it can plan a route several rounds ahead. It is emphatically not given
the seed the \archipelago{} was generated from: that seed determines where the \treasure{} is, and a
\ship{} holding it could simply re-run generation and sail there. The \chart{} must therefore derive
two independent streams from the game seed, and only the weather one ever leaves it. This is easy to
get wrong and impossible to notice from the outside, so it is a requirement rather than a
suggestion.

\textbf{The roster is complete and public.} Who is playing is known to everybody from the start.
Where they are is another matter entirely, and \texttt{Welcome} carries no position.

\textbf{There is no \texttt{mode} field, and there must not be one.} The \MD{} requires that no
message a \ship{} receives name the schedule it is being run under, and a field saying so would
hand it the answer outright. It would also be read: tell a competitor which mode it is in
and somebody will branch on it, and gain in Sequential what it could not gain in Simultaneous.
The \cartographer{} is told the mode instead (\cref{proto:cartographer}), because no \ship{} can
receive a message addressed to it --- which is the whole reason there are four message sets
(\cref{proto:foursets}).

\newpage
\section{\clues{}}
\label{proto:clue}

A \clue{} is a predicate over cells. A \clue{} that came from an \island{} is true of the
\treasure{}'s cell; a \clue{} that came from another \ship{} may be anything at all. The \MD{}
explains what \clues{} are for and where they come from. This section fixes what one \emph{is}.

\subsection{The type}

\begin{lstlisting}
#[non_exhaustive]
#[derive(Clone, Copy, PartialEq, Eq, Hash, PartialOrd, Ord, Debug)]
pub enum Clue {
    /// No further than `radius` from `centre`.
    Within { centre: Cell, radius: u32 },
    /// Further than `radius` from `centre`.
    Beyond { centre: Cell, radius: u32 },
    /// The coordinate on `axis` satisfies `bound`.
    HalfPlane { axis: Axis, bound: Bound },
    /// Strictly closer to `a` than to `b`.
    Nearer { a: &Cell, b: Cell },
}

#[derive(Copy, Clone, PartialEq, Eq, Hash, PartialOrd, Ord, Debug)]
pub enum Axis { Q, R, S }

pub enum Bound { AtLeast(i32), AtMost(i32) }

impl Clue {
    /// True when this clue is consistent with the treasure lying in `cell`.
    pub fn holds(self, cell: Cell) -> bool {
        match self {
            Clue::Within { centre, radius } => cell.distance(centre) <= radius,
            Clue::Beyond { centre, radius } => cell.distance(centre) >= radius,
            Clue::HalfPlane { axis, bound } => {
                let v = match axis {
                    Axis::Q => cell.q,
                    Axis::R => cell.r,
                    Axis::S => cell.s(),
                };
                match bound {
                    Bound::AtLeast(k) => v >= k,
                    Bound::AtMost(k) => v <= k,
                }
            }
            Clue::Nearer { a, b } => cell.distance(a) < cell.distance(b),
        }
    }
}
\end{lstlisting}

\subsection{Why it looks like this}

\textbf{Four variants, and they cut the sea in different shapes.} \texttt{Within} and
\texttt{Beyond} give discs and their complements, and a pair of them gives an annulus.
\texttt{HalfPlane} gives a straight cut along one of the three hex axes. \texttt{Nearer} gives a
perpendicular bisector. Intersecting four \clues{} of the same shape is arithmetic; intersecting
four of these is a puzzle, which is the point.

\textbf{There is no \texttt{Not(Box<Clue>)}.} Recursion would buy expressiveness the game does not
need, at the cost of an allocation and a recursive \texttt{holds}. The negations that matter already
exist as flat variants: \texttt{Beyond} is the complement of \texttt{Within}, and \texttt{Bound}
gives both directions of a half-plane.

\textbf{\texttt{holds} takes a cell and nothing else.} It is pure, it is total, and it cannot
panic: every arm is a comparison between integers. None of the four variants needs to consult the
map. Should the \WG{} ever add one that does --- ``not adjacent to land'' is the obvious candidate
--- the signature grows a \texttt{\&Map} at that point and not before.

\textbf{\texttt{Nearer} names cells, not \islands.} \islands sit on cells that everybody knows, so
\texttt{Cell} is strictly more general and saves a lookup. An \island{} is still free to announce
the \clue{} as ``closer to the northern \island{} than to the southern one''; that is presentation,
and it is the
\island's business.

\textbf{The derives are load-bearing.} \texttt{Ord} and \texttt{Hash} exist so that a collection of
\clues{} can be put in a deterministic order before anything iterates it, which the determinism
rules require. \texttt{Copy} exists because a \texttt{Clue} is four words and copying it is cheaper
than reasoning about who owns it.

\textbf{The variant is part of the \clue, not an implementation detail.} When an \island{} weakens
a \clue{} it answers with the same variant it was given --- a weakened \texttt{Nearer} is a
\texttt{Nearer}. That is why weakening is a matter of moving one field, and
why \texttt{Clue} needs no separate notion of strength.

\textbf{\texttt{\#[non\_exhaustive]} is deliberate}, and it is aimed squarely at \ships. A \ship{}
should ask a \clue{} whether it holds, never take it apart to see what kind it is. The attribute
makes that the path of least resistance: a \ship{} that matches on variants must write a \texttt{\_}
arm, and a \ship{} that leans on \texttt{holds} keeps working unchanged when a fifth variant
appears.

\subsection{What a \ship{} does with them}

The cells still consistent with a collection of \clues{} are the water cells that satisfy all of
them:

\begin{lstlisting}
let candidates: Vec<Cell> = map
    .water_cells()
    .filter(|&c| clues.iter().all(|clue| clue.holds(c)))
    .collect();
\end{lstlisting}

A \ship{} that has read the \MD{} can do better than that in one line. The \treasure{} is never on
a rock and never adjacent to land, so \texttt{clear\_cells} is the right iterator and the cells
touching land can go too --- which on a crowded map throws away a good fraction of the sea before
a single \clue{} is evaluated. The snippet above is correct; it is not the best you can do.

This requires a \ship{} to know the map, and it does: \textbf{the geography is public}. The kind of
every cell --- clear water, rock, \island{} footprint, cliff --- is delivered with
\texttt{Welcome} and never changes. A \ship{} needs it to plan a route in any case. What is hidden
is where the rivals are, and which single cell holds the \treasure{}.

Note also that the snippet above is the floor and not the ceiling. A \ship{} that does only this has
nothing left to say the moment \texttt{candidates} comes back empty --- which it will, the first
time a rival lies to it.

\newpage
\section{Rounds}
\label{proto:rounds}

\subsection{The exchange}

Each round is one exchange with each \ship{} still afloat:

\begin{lstlisting}
/// Admiralty -> Ship. The round is open; answer before the deadline.
AdmiraltyToShip::RoundOpen {
    round: Round,
    /// Milliseconds from receipt, not an absolute time: clocks are not synchronised.
    deadline_ms: u32,
    view: ShipView,
    token: TurnToken,
}

/// Ship -> Admiralty. Exactly one action, plus anything to say.
ShipToAdmiralty::Submit {
    round: Round,
    token: TurnToken,
    action: Action,
    outbox: Vec<Letter>,
}

/// Admiralty -> Ship, after the round is resolved.
AdmiraltyToShip::RoundResult {
    round: Round,
    view: ShipView,
    events: Vec<Event>,
    inbox: Vec<Delivered>,
}

/// Admiralty -> Ship. The last message a Ship ever receives.
AdmiraltyToShip::GameOver {
    round: Round,
    reason: EndReason,
    /// What this Ship is told about the result. A Ship still playing when the
    /// game ends is told everything; one that sank in round twelve cannot be,
    /// because there was nothing to tell yet.
    outcome: Outcome,
}

#[derive(Copy, Clone, PartialEq, Eq, Debug)]
pub enum EndReason {
    /// This Ship sank. It is out; the game continues without it.
    Sunk,
    /// Somebody declared the Treasure correctly.
    Declared,
    /// The round limit was reached.
    RoundCap,
}
\end{lstlisting}

\textbf{There is no variant for ``everybody sank''}, and adding one would be a mistake. The \MD{}
ends the game when no \ship{} is left afloat, but by then all six have already been told
\texttt{Sunk} at the moment they went down, and each receives exactly one \texttt{GameOver}. There
is nobody left to deliver a seventh message to. That ending is visible to the \chart{}, through
\texttt{outcome}, and to the \cartographer{} --- and to nobody else, because there is no longer
anybody else.

\texttt{deadline\_ms} is a duration and not a moment, so that a \ship{} never has to read a clock
to know how long it has. It is measured by the \admiralty{} from the point of sending.

\textbf{\texttt{GameOver} is what lets a \ship{} thread stop.} Without it a sunk \ship{} would
simply stop receiving \texttt{RoundOpen} and would block for ever on an endpoint nobody will
write to again --- which is the one failure this project spends a whole section forbidding. The
\MD{} requires that a sunk \ship{} be told its game is over; this is the message that tells it.
Every \ship{} receives exactly one, and receives nothing afterwards.

A \ship{} that is told \texttt{Sunk} is out while the game runs on. It should drop its endpoints
and return; it must not keep submitting, and anything it does send is discarded.

\texttt{TurnToken} is issued with the round and consumed by the submission. A \ship{} built around
it cannot express the idea of acting out of turn: there is nothing to pass. Note what this does and
does not achieve --- it keeps honest code correct, and it stops nothing at all in a \ship{} that
simply stores an old token and sends it again, which is why the \admiralty{} checks the round
number as well.

\subsection{The order of business}

For each round the \admiralty:

\begin{enumerate}
    \item sends \texttt{RoundOpen} to every \ship{} still afloat, each with its
      own \texttt{ShipView} obtained from the \chart;
    \item collects submissions until the deadline;
    \item for each submission, in the standardised order (\fullref{proto:conflicts}), asks the
      \chart{} to validate and then \emph{immediately} to apply, answering a refusal with
      \texttt{Rejected};
    \item calls \texttt{advance\_round}, the \chart's end-of-round step;
    \item sends every \ship{} a \texttt{RoundResult} carrying its own view and its own inbox;
    \item sends the \cartographer{} the complete state.
\end{enumerate}

That is the order in Simultaneous mode.

Step 3 says \emph{immediately} for a reason. A \texttt{ValidatedAction} carries no round and no
state version, so nothing in the type system stops an \admiralty{} from validating all six
submissions and then applying all six --- and that \admiralty{} will put two \ships{} in one cell,
because the second move was judged against a sea the first had not yet moved in. The \MD{} judges
legality at the moment an action is applied; validate and apply must therefore be adjacent, with
nothing in between.

A \ship{} that has not answered by the deadline is recorded as having played \texttt{Hold}. It is
never removed from the game for being slow. A thread that thinks too hard, blocks, or dies must not
be able to stop the game for everybody else.

\subsection{Two modes, one \ship}
\label{proto:modes}

The \admiralty{} runs in one of two modes:

\begin{lstlisting}
#[derive(Copy, Clone, PartialEq, Eq, Debug)]
pub enum TurnMode { Sequential, Simultaneous }
\end{lstlisting}

In \textbf{Sequential} the round is opened for one \ship. In \textbf{Simultaneous} it is opened
for all of them.

\textbf{No message a \ship{} receives names the mode, and none may be added that does.} The mode is
a scheduling decision belonging to the \admiralty. The
\cartographer{} is told, because it draws the game and must be able to say what it is looking at.

This works because \texttt{RoundOpen} does not say ``it is your turn''. It says: round $R$ is open,
here is what you can see, answer within this long. Whether one \ship{} received that message or
all six did is not visible to any of them. The \ship{} code is identical.

Four rules keep it that way. Each is load-bearing; drop any one and the modes stop being
interchangeable:

\begin{itemize}
    \item no message a \ship{} can receive carries the mode, and none may be added that does;
    \item \texttt{RoundOpen} carries the round number and a deadline, and never implies
      exclusivity;
    \item every \texttt{Submit} names its round, and a submission naming a closed round is
      discarded;
    \item results are per-\ship{} and keyed by round, so a \ship{} can always tell which round it is
      being told about.
\end{itemize}

The \MD{} is careful about how far this goes, and so should you be. The modes are not
statistically identical --- a \ship{} in Sequential is never refused a cell a rival took in the
same round --- so a \ship{} that collected its rejections could form an opinion. The requirement
is that nothing tells it and that no \ship{} acts on a guess. Do not design a message that would
make the guess unnecessary.

Sequential exists because it is reproducible and easy to follow: it is the mode for testing and for
examinations. Simultaneous exists because a round takes the same time however many \ships have
stalled, which is what a game of six thinking players needs.

\subsection{Conflicts}
\label{proto:conflicts}

Simultaneous rounds allow what Sequential cannot: two \ships claiming the last berth, firing on
each other at once, or declaring the same cell in the same round.

\textbf{The \WG{} must standardise a total, deterministic order over the submissions of a round,
and the outcome of every class of conflict.} It is not enough for the order to be deterministic
within one implementation: it must be the same order in every implementation, or the same
submissions will produce different games and neither replay nor any comparison between
implementations will mean anything. The same order also fixes the sequence in which Sequential
asks (\cref{proto:modes}), so this one decision governs both modes.

The order must be a function of information every implementation has. Arrival order is not: it
depends on which thread the operating system happened to wake first, which is precisely the
nondeterminism this project forbids. Plausible bases are the \ship{} identifier, or a permutation
derived from the seed and the round number --- the latter being fairer, since it does not advantage
a low identifier for the whole game.

Conflicts the \WG{} must give an answer for, at least:

\begin{itemize}
    \item two \ships moving into the same cell, which only one may have: cells are exclusive;
    \item a \ship{} moving into the cell another \ship{} is leaving in the same round, which
      succeeds or fails purely according to which of the two is applied first;
    \item more \ships mooring than there are berths;
    \item two \ships bombarding each other, when the first shot would sink the second: does the
      second shot happen;
    \item two \ships landing a killing blow on the same target, which decides who inherits the
      victim's \clues{};
    \item a \ship{} moving out of a cell that another \ship{} is firing into.
\end{itemize}

Declarations are not on that list, and it is worth knowing why: \texttt{Declare} names the \ship's
own cell and cells are exclusive, so two \ships can never both be right in the same round.

The killing blow deserves a moment. Two \ships have hunted the same rival for forty rounds and both
fire the shot that sinks it in the same round; only one of them inherits everything that \ship{}
knew. Which one is decided entirely by the conflict rule. It is the most consequential sentence the
\WG{} will write, and it will be written in a meeting about tie-breaking.

\newpage
\section{The \chart's interface}
\label{proto:chartapi}

The \chart{} owns no channel and sends nothing. The \admiralty{} \emph{calls} it, synchronously,
and this is the only interface between the two --- which makes it as much a part of what two
groups must agree on as any message, and the reason it is in this document rather than hidden
inside whoever wrote the \chart.

\begin{lstlisting}
impl Chart {
    pub fn generate(seed: Seed, params: &Params) -> Result<Chart, GenError>;

    /// Decide whether an action is legal. The only way to obtain a ValidatedAction.
    pub fn validate(&self, ship: ShipId, action: &Action)
        -> Result<ValidatedAction, Rejection>;

    /// Apply an already validated action. Cannot fail.
    pub fn apply(&mut self, action: ValidatedAction) -> Vec<Event>;

    /// End of round: the round number advances and the Islands take their turn.
    pub fn advance_round(&mut self) -> Vec<Event>;

    /// What a Ship is permitted to see.
    pub fn view_for(&self, ship: ShipId) -> ShipView;

    /// What a spectator sees: everything.
    pub fn spectator_view(&self) -> HashMap<ShipId, ShipView>;

    pub fn outcome(&self) -> Option<Outcome>;
}
\end{lstlisting}

\subsection{Why validate and apply are two functions}

\texttt{ValidatedAction} has no public constructor: it cannot be built outside the \chart's own
crate, and \texttt{validate} is the only thing that returns one. \texttt{apply} accepts nothing
else. It is therefore not merely discouraged but \emph{impossible} to apply an unchecked action to
the game state, and the compiler enforces it rather than a code review.

This is worth copying elsewhere in your own code. A type whose values a caller cannot construct is
how you make an invariant unforgeable in Rust: the invariant lives in the one function that mints
the value, and everything downstream may simply assume it.

\texttt{apply} cannot fail, and its signature says so. Everything that might have gone wrong was
decided by \texttt{validate}. If you find yourself wanting to return an error from \texttt{apply},
the check belongs in \texttt{validate}.

\textbf{Note the one thing the type does not buy you.} A \texttt{ValidatedAction} says the action
was legal when it was minted; it does not say it is legal now, because it carries no round and no
state version. Hold one across another \texttt{apply} and it may have gone stale --- the cell it
names may have filled, its target may have sunk. The \MD{} judges legality at the moment of
application, so \texttt{validate} and \texttt{apply} must be adjacent calls with nothing between
them (\cref{proto:rounds}). The type stops you applying an \emph{unchecked} action. Keeping the
check fresh is the \admiralty's job.

\subsection{What the \admiralty{} may rely on}

\begin{itemize}
    \item \texttt{validate} never mutates, so a rejected action leaves the state untouched and the
      \admiralty{} may call it speculatively;
    \item \texttt{apply} returns the \texttt{Event}s its action produced, and those events are the
      only record of what happened;
    \item \texttt{advance\_round} does three things in this fixed order: increment the round
      number, so the wind takes its new value; give every \island{} its turn, in ascending order
      of identifier, so that a \ship{} moored and waiting has waited one round more; determine
      whether the game has ended --- because somebody declared correctly, because the round cap
      was reached, or because no \ship{} is left afloat. The order is fixed and the traversal is
      sorted because both affect the outcome and both must be identical in every implementation;
    \item \texttt{view\_for} is pure and may be called as often as the \admiralty{} likes;
    \item there is no sinking step. A \ship{} sinks inside the \texttt{apply} that empties its
      hull, never at the end of a round (see the \MD).
\end{itemize}

Everything the \chart{} does with these calls --- how a world is generated, what makes an action
legal, what a view may contain --- is in the \MD{} with the rest of the rules.

\newpage
\section{Messages between \ships}
\label{proto:mail}

\subsection{How they travel}

A \ship{} may send messages to other \ships. They are carried in the \texttt{outbox} of a
submission and delivered in the \texttt{inbox} of a result.

\begin{lstlisting}
/// What a Ship sends.
pub struct Letter {
    pub to: ShipId,
    /// Need not be clues the sender holds. Lying is the point.
    pub clues: Vec<Clue>,
    pub text: String,
}

/// What a Ship receives. The Admiralty adds `from`.
pub struct Delivered {
    pub round: Round,
    pub clues: Vec<Clue>,
    pub text: String,
}
\end{lstlisting}

\textbf{The sender is stamped by the \admiralty{} and cannot be forged.} A \ship{} chooses what it
says and to whom; who it is comes from the endpoint the \texttt{Letter} arrived on, and is not a
field the sender can set. This is the one thing about a message that is guaranteed. Everything in
it may be a lie; the name on it may not.

\clues{} travel as \texttt{Clue} values rather than as prose, because the receiver has to be able to
evaluate them (see the \MD). \texttt{text} is for everything else --- promises, threats,
negotiations --- and nothing in the system reads it but the \ship{} it is addressed to.

A message sent in round $R$ arrives with the results of round $R$, and can therefore be read before
choosing an action for round $R+1$. Coordination costs one round of delay.

Sending is not an action: a \ship{} acts and speaks in the same round. The limits are \textbf{5}
letters per submission, and at most \textbf{4} \clues{} and \textbf{256} characters in each.
Something must bound it, or six \ships over two hundred rounds will produce logs too large to
read.

A submission that exceeds any of these is \textbf{rejected, not quietly trimmed}. Silently dropping
the fourth letter would leave a \ship{} believing it had said something it had not, and its author
hunting a bug in the wrong component. Tell it. Note that the whole submission goes, action
included: the \ship{} is answered with \texttt{Rejection::TooMuchMail} and recorded as having
played \texttt{Hold}. Overrunning the cap costs a round.

The \admiralty{} makes a reasonable effort to deliver and no promise. A letter addressed to a
\ship{} that does not exist, to one that has sunk, or to the sender itself is dropped. There is no
acknowledgement, no delivery report, and no ordering guarantee between letters from different
senders.

\subsection{Why it is relayed and not direct}

Because two independently written student components that wait on each other will wait forever. A
\ship{} blocked on a peer during its own round stalls nothing but itself --- the \admiralty{} will
time it out --- but a \ship{} that has staked its turn on an answer that never comes has wasted the
round, and at the fair it will do so in front of an audience.

Uncooperativeness is also a legitimate strategy here. A \ship{} that depends on being answered is a
\ship{} its rivals can freeze simply by saying nothing, and the relay is what keeps that from being
a winning move: the mail never blocks anybody, because nobody ever waits for it.

\subsection{The two rules}

\begin{description}
    \item[A message may never be a precondition for the legality of an action.] The \chart{}
      decides legality from its own state alone. It never sees a message and does not know any were
      exchanged. Anything agreed between \ships must be re-stated as an ordinary action to have any
      effect.
    \item[A \ship{} must remain fully playable if every other \ship{} ignores it forever.] This is
      tested: a \ship{} is run in a game where no peer ever replies, and must complete it within the
      normal deadlines. A \ship{} that fails does not pass integration.
\end{description}

The second rule is the only mechanically checkable proxy for the first, which is why it is a test
and not advice.

\subsection{What this is for}

The consequence is the point of the whole arrangement: \textbf{what is negotiated is private and
cannot be verified; what is done is public and is validated.}

You can tell another \ship{} that an \island{} swore the \treasure{} lies east of a given line, and
it cannot check. You can agree to attack on round forty and simply not turn up. The \admiralty{}
sees who fired; it never knew a promise existed, and cannot be appealed to. Deception is not an
exploit here. It is the mechanic.

\newpage
\section{The \cartographer}
\label{proto:cartographer}

\subsection{Watching}

The \cartographer{} receives the complete state, unfiltered, every round:

\begin{lstlisting}
AdmiraltyToCartographer::State {
    round: Round,
    state: SpectatorView,
    events: Vec<Event>,
    /// Which schedule is running. No Ship is ever told this.
    mode: TurnMode,
}
\end{lstlisting}

\texttt{SpectatorView} is its own type and is not \texttt{ShipView} with the filtering disabled.
The whole point is that the complete state is not expressible in anything a \ship{} can receive.

Relayed messages are visible to the \cartographer, since somebody has to be able to see what was
said when the game is over and two groups disagree about what happened.

\subsection{Controlling}

The \cartographer{} may also drive the game:

\begin{lstlisting}
CartographerToAdmiralty::
    Pause,
    Resume,
    Step,                                  // advance exactly one round, then pause
    SetMode { mode: TurnMode },
    SetDeadline { deadline_ms: u32 },
    ExpireTurn { ship: ShipId },           // treat as Hold immediately
    Dump,                                  // full state, out of band
\end{lstlisting}

\texttt{Step} with \texttt{Pause} is what makes a game examinable: a round at a time, with time to
ask why a \ship{} chose what it chose.

\texttt{SetDeadline} changes the deadline for the rounds that follow. Note that it leaves the
\texttt{round\_deadline\_ms} every \ship{} was handed in \texttt{Welcome} out of date, which is
why the authoritative figure is the \texttt{deadline\_ms} in \texttt{RoundOpen}: that one arrives
every round and is always current. A \ship{} that budgets its thinking from \texttt{Params} rather
than from the round it is answering will be timed out by a \cartographer{} tightening the screw
during a demonstration.

\textbf{These commands are privileged, and the four message sets are what enforces it.} They live
only in \texttt{CartographerToAdmiralty}. A \ship{} is handed its own pair of endpoints
and cannot acquire others (\cref{proto:connect}), so there is no way for a competitor to express
\texttt{Pause} or \texttt{ExpireTurn} at all --- it is a compile error, not a check the
\admiralty{} has to remember to write. This is the one place the project needs authorisation, and
it is worth seeing that it was obtained by choosing the types rather than by adding a rule.

The \cartographer{} is optional. A game with nobody watching runs perfectly well.

\newpage
\section{Errors and hostile input}
\label{proto:errors}

\subsection{Rejections}

A submission the \chart{} refuses is answered, not ignored:

\begin{lstlisting}
AdmiraltyToShip::Rejected { round: Round, reason: Rejection }
\end{lstlisting}

\texttt{Rejection} is an enumeration, not a string. Sketch, to be completed by the \WG:

\begin{lstlisting}
#[non_exhaustive]
#[derive(Clone, PartialEq, Eq, Debug)]
pub enum Rejection {
    WrongRound { open: Round }, StaleToken,
    NotANeighbour, OutsideArchipelago, CellIsLand,
    CellOccupied { by: ShipId },
    NotAdjacentToIsland, NoBerth, IslandRefused, NotMoored,
    OutOfRange, NotVisible,
    NoSuchShip, AlreadyAllied, NoSuchProposal, NotAllied,
    TargetHasSunk,
    TooMuchMail,
}
\end{lstlisting}

\texttt{IslandRefused} covers both ways an \island{} can turn a \ship{} down: refusing it a berth
when it plays \texttt{Moor}, and refusing to answer when it plays \texttt{AskClue}. One variant
for one idea.

Whoever wrote the \ship{} cannot ask you what went wrong, and will be debugging late at night
against an \admiralty{} written by somebody they have never met. A reason they can match on is
worth a great deal more than a sentence they have to read.

A rejected action leaves the state untouched, and the round proceeds as though the \ship{} had
played \texttt{Hold}. The \ship{} is not removed from the game for submitting nonsense, however
much of it it submits.

\subsection{Hostile input}

Every value arriving at the \admiralty{} was written by a competitor.

That the messages are typed Rust values rather than bytes removes one class of problem and none of
the others. Nobody can hand you a truncated struct or a length prefix that lies. Everybody can hand
you a \texttt{Cell} outside the sea, a \texttt{ShipId} that names nobody, an action for a round
that ended an hour ago, an outbox with a hundred thousand letters in it, or a \texttt{String} the
size of the heap.

\textbf{The \admiralty{} must not panic, abort or exit on any message, however absurd.} Somebody
will send each of the above, some of them on purpose. Treat every value crossing a channel as an
assertion by an adversary and check it before use --- including the ones that ``obviously'' cannot
be wrong.

A misbehaving \ship{} may be silenced. The game may not stop.

\newpage
\section{Logging}
\label{proto:logging}

Every component logs the messages it sends and receives, through \texttt{log} or \texttt{tracing}
(\url{https://docs.rs/log/latest/log/}, \url{https://docs.rs/tracing/latest/tracing/}). The \WG{}
standardises the format and the event targets, so that logs from components written by different
groups can be read side by side. That is the entire value of standardising it: a log you cannot
line up against somebody else's answers no question worth asking.

The \admiralty{} additionally records, for the whole game: the seed, the parameters, the turn mode,
and every submission it accepted with its round number. That record plus the seed reproduces the
game exactly (see the \MD{} on determinism), which is what makes a finished game replayable and a
disagreement settleable.

Relayed messages are logged too. Their contents prove nothing --- they may be lies --- but their
existence and their timing can at least be established afterwards.

\newpage
\section{What the \WG{} must settle}
\label{proto:todo}

The \MD{} lists the decisions the \WG{} owes for the game's \emph{rules}. What follows is the same
list for its \emph{messages}, and each item is open because it is a decision rather than an
oversight:

\begin{enumerate}
    \item \textbf{The conflict order and every conflict outcome} (\fullref{proto:conflicts}). The
      most important item on this list, and the one with the largest consequences.
    \item \textbf{The \texttt{Action} enumeration.} The \MD{} lists the actions and what they mean.
      \texttt{Move} is settled --- it carries the single neighbouring \texttt{Cell} to sail to ---
      but the rest is not: how an \island{} is named, and what an \island's answer to
      \texttt{AskClue} looks like when it is neither a \clue{} nor a refusal. Note that
      \texttt{AcceptAlliance} has to name a proposer: a \ship{} may hold one outgoing proposal but
      any number of incoming ones (see the \MD).
    \item \textbf{The \texttt{Event} enumeration.} What a \ship{} is told happened, as distinct from
      what it can see. Take care: an event is as capable of leaking information as a view is, and it
      is easier to forget. What a bombardment reveals is already settled in the \MD{} --- the
      victim learns who fired, the attacker learns whether it hit --- but the shape of the events
      that carry it is not. A \ship{} striking a rock needs an event too, and so does a \ship{}
      sinking on one with nobody to inherit its \clues{}. \textbf{Three things do not follow line
      of sight}, and the \MD{} fixes all three: an alliance formed or broken, which goes to all
      six \ships; a declaration, which goes to all six and carries the cell that was named; and
      the name of whoever fires a shot, which reaches the victim and every onlooker who can see
      the victim's cell, even when the attacker is nowhere any of them could see. Everything else
      is filtered by what the \ship{} can see, and no fourth exception may be invented here.
    \item \textbf{\texttt{ShipView} and \texttt{SpectatorView}} in detail. What may be in them is
      settled in the \MD; their shape is not.
    \item \textbf{The \island{} trait}, likewise. Two of its methods are not optional: the \MD{}
      obliges an \island{} to be able to say that it has nothing further to offer a given \ship,
      and to say how much longer that \ship{} must wait. Both have to exist in the trait, or an
      \island{} cannot keep a promise the \MD{} makes on its behalf. The answer type for
      \texttt{AskClue} belongs here as well, and must be able to express all three legal outcomes:
      a \clue, a refusal, and ``not yet, come back in $n$ rounds''. A \clue{} it returns is always
      of the same variant as the one it was given.
    \item \textbf{\texttt{Outcome}}, carried by \texttt{GameOver} (\fullref{proto:rounds}) and
      returned by \texttt{Chart::outcome}. The \MD{} fixes what is ranked --- how many of the
      world's \clues{} a \ship{} holds at full strength, ties by identifier --- so what is open is
      the shape: whether a \ship{} is told the whole ordering or only its own place, and how much
      of either a \ship{} that sank in round twelve can be told, given that the game had not
      finished when its \texttt{GameOver} was sent. Note two things the \chart{} must carry for
      it: a sunk \ship's count, recorded as it sinks, because its \clues{} pass to the killer a
      moment later; and the fact that a game ended with nobody afloat, which no \texttt{EndReason}
      expresses.
    \item \textbf{\texttt{GenError}}, returned by \texttt{generate}. The \MD{} gives it at least
      two causes --- parameters refused outright, and rejection sampling that ran out of attempts
      --- and a harness that cannot tell them apart cannot say whether to fix the parameters or
      try another seed.
    \item \textbf{Deadline behaviour} when a submission arrives in the same instant it expires.
    \item \textbf{Channel capacity}, and what the \admiralty{} does when a \ship's inbound channel
      is full.
\end{enumerate}

One piece of advice on how to run this. The productive way to find the holes in a protocol is to
try to abuse it. Take a message and ask what happens if it arrives twice, arrives late, arrives
before \texttt{Welcome}, or arrives from somebody who should not be able to send it. Draw the
states a \ship{} can be in --- sailing, moored, sunk --- and look for the transitions nobody has
named. The problems worth finding are almost never in a single message; they are in two features
that were each specified alone.

\end{document}
